\scanpage{57}
% archive sheet number: 56
\deleted{Fixons \unsure{une|un} $v$ de référence (i.e. une topologie sur $k$, \unsure{donnant|donnée} \unsure{isom} $k\unsure{\mathrel{\text{---}}|\longrightarrow}\mathbb C$ mod conj. compl.).}

Or $g_\sigma$ est connu mod mult. à droite par $G(\mathbb Z)$, d'autre part $g_{\sigma\,\mathrm{fini}}$ est (avec la = indétermination) le système des $\sigma_l$ galoisiens. Donc on peut prendre
\[
g_\sigma=((\deleted{\unsure{g_l}}\sigma_l)_{l\ \mathrm{fini}},\sigma_\infty)
\]
\uline{où maintenant} $\sigma_\infty$ \uline{est un élément bien déterminé de} $G(\deleted{\unsure{\mathbb C}}k)$ !!!

On obtient une application canonique
\diagram{\fbox{\begin{tikzcd}[ampersand replacement=\&]
\operatorname{Gal}(k/k_0)\arrow[r,"{\varphi_{v\infty}}"]\&G_{\deleted{\unsure{M}}}(\deleted{\unsure{\mathbb C}}k)
\end{tikzcd}}$\subset\mathrm{Gl}(M_k(\infty))$}
\margin{NB. $G$ peut ici être considéré comme groupe défini sur $k_0$ (indépendant de $v$ !)}
\deleted{et} telle qu'on ait
\diagram{\begin{tikzcd}[ampersand replacement=\&,row sep=small,column sep=small]
\dot{\sigma}M_v^{\mathbb Z}\arrow[r,equal,end anchor=north west]\&\left\{\begin{aligned}
M_{\sigma v}^{\mathbb Z}&=\sigma_*M_v^{\mathbb Z}\quad\text{pour }\sigma\in\operatorname{Gal}(k/k_0)\\
\sigma_*&=((\sigma_l)_{l\ \mathrm{fini}},\varphi_v(\sigma)).
\end{aligned}\right.
\end{tikzcd}}
On vérifie aussitôt
\[
\varphi_v(\tau\sigma)={}^{\tau}\varphi_v(\sigma)^{\deleted{\unsure{\tau}}}\varphi_v(\tau)
\]
i.e. posant $\psi_v=\varphi_v^{-1}$
\[
\psi_v(\tau\sigma)=\psi_v(\tau)\,{}^{\tau}\psi_v(\sigma)^{\deleted{\unsure{\tau}}}
\]
